\begin{helper}
Let \(K\) be a finite field of order \(q\), let \(t\ge2\), let
\(A,k,w,\Delta\) be nonnegative integers, and let \(C\in\mathbb R\).  Assume
the projective parameter upper bounds
\[
  \frac{q^{t+1}-1}{q-1}\le 2q^t,\qquad
  \frac{q^t-1}{q-1}\le 2q^{t-1},
\]
and assume
\[
  4q^{2t-1}\le\Delta,\qquad 10128t\le A.
\]
Assume also, in real-number inequalities, that
\[
  A\le 4q^{t-1},\qquad A\le C,\qquad 1\le\log q,
\]
\[
  w\le Aq\log q,\qquad Cq(\log q)^2\le k,
\]
and that the concrete shrink threshold holds:
\[
  \left(1-\frac{1}{32tq}\right)^{w/(t+1)-1}
  < \frac{1}{4q^t},
\]
where the exponent \(w/(t+1)-1\) is the truncated natural-number exponent.
Then, for the concrete marking rule
\[
  \noderef{StarProductConcreteMarked}
  \quad\text{with}\quad
  \ell(p,r)=\noderef{StarProductLayerChoice}(p,r),
\]
one has
\[
  \fwi_k(D^*(G(t,K)))
  \le 2^k\Delta^w(Aq^t)^{k-w}.
\]
\end{helper}
\begin{proof}
We verify the hypotheses of the tree-counting bridge
\(\noderef{StarProductMarkedTreeCountingBridge}\) for
\[
  G=G(t,K),\qquad
  h=Aq^t,
\]
and for the concrete marking rule displayed in the statement.

First, the finite-field parameter estimate
\(\noderef{StarProductPolarityParameterBounds}\) gives \(q\ge2\), and hence
\(q\ge1\).  Since \(h=Aq^t\), the inequality \(h\le Aq^t\) required by
\(\noderef{StarProductHLeDelta}\) is an equality.  Together with the
assumptions \(A\le4q^{t-1}\) and \(4q^{2t-1}\le\Delta\), that node gives
\[
  h=Aq^t\le\Delta.
\]

Next, applying \(\noderef{StarProductWLeK}\) with the real parameters
\(A\) and \(C\) gives
\[
  w\le k.
\]
Indeed its hypotheses are exactly \(A\ge0\), which is automatic because
\(A\) is a natural number, together with \(A\le C\), \(1\le\log q\),
\(w\le Aq\log q\), and \(Cq(\log q)^2\le k\).

For the path hypothesis in
\(\noderef{StarProductMarkedTreeCountingBridge}\), fix a forward independent
\(k\)-tuple \(v\) in \(D^*(G(t,K))\).  The shrink-threshold assumption is
precisely the hypothesis of
\(\noderef{StarProductConcretePathWeightBound}\), so that node gives
\[
  \bigl|\operatorname{sig}(v)\bigr|\le w
\]
for the concrete marked-tuple signature.

For the all-children hypothesis, the assumed upper bounds on the projective
vertex and degree parameters and the inequality \(4q^{2t-1}\le\Delta\) are
exactly the hypotheses of \(\noderef{StarProductConcreteAllChildrenBound}\).
Thus every legal forward-independent prefix of length \(m<k\) has at most
\(\Delta\) legal children.

For the marked-children hypothesis, apply
\(\noderef{StarProductConcreteMarkedChildrenBound}\) with \(h=Aq^t\).  Its
condition \(Aq^t\le h\) is again equality, and the remaining numerical
hypothesis is \(10128t\le A\), which is assumed.  Therefore every legal
prefix of length \(m<k\) has at most \(h=Aq^t\) legal children marked by the
concrete rule.

All hypotheses of \(\noderef{StarProductMarkedTreeCountingBridge}\) are now
verified: \(w\le k\), \(h\le\Delta\), every full path has at most \(w\)
unmarked steps, every prefix has at most \(\Delta\) legal children, and every
prefix has at most \(h\) marked legal children.  Applying the bridge yields
\[
  \fwi_k(D^*(G(t,K)))\le 2^k\Delta^w h^{k-w}.
\]
Substituting \(h=Aq^t\) gives exactly
\[
  \fwi_k(D^*(G(t,K)))\le 2^k\Delta^w(Aq^t)^{k-w}.
\]
\end{proof}
